\section[Penggunaan PHP dari Command Line (CLI)]{Penggunaan PHP dari Command Line (CLI)}

PHP tidak hanya untuk web --- PHP juga dapat dijalankan dari \textbf{command line} (CLI) untuk skrip otomasi, cron job, migrasi database, atau pengujian. Perintah \code{php namafile.php} menjalankan skrip; \code{php -r "echo 'test';"} menjalankan kode langsung. Variabel \code{\$argv} berisi argumen CLI, dan \code{\$argc} jumlahnya. PHP CLI berguna untuk tugas backend yang tidak memerlukan browser --- misalnya import data massal atau kirim email terjadwal \cite{phpManual}.

\begin{webcode}[language=PHP, caption={Contoh skrip PHP CLI}]
<?php
// file: migrate.php
// Jalankan: php migrate.php create_table

if ($argc < 2) {
  echo "Penggunaan: php migrate.php <perintah>\n";
  exit(1);
}

$command = $argv[1];
echo "Menjalankan migrasi: {$command}\n";

// Contoh: buat tabel jika belum ada
$conn = new mysqli('localhost', 'root', '', 'toko_online');
$sql = "CREATE TABLE IF NOT EXISTS log_migrasi (
  id INT AUTO_INCREMENT PRIMARY KEY,
  nama VARCHAR(100), waktu DATETIME DEFAULT CURRENT_TIMESTAMP
)";
$conn->query($sql);
echo "Migrasi selesai.\n";
?>
\end{webcode}

